\section{Software System Design}

This section describes the software system surrounding the segmentation model. NephroVision is implemented as a browser-accessible application that ingests CT volumes, executes the segmentation pipeline, and presents results through an interactive web interface. The system is designed as an academic prototype for decision-support use and is not deployed in any clinical environment.

\subsection{System Architecture Overview}

NephroVision follows a client-server architecture comprising three main components:

\begin{itemize}
    \item \textbf{Frontend:} A browser-based user interface implemented using standard web technologies (HTML, CSS, JavaScript). The frontend handles file upload, displays segmentation results, and renders interactive 3D visualizations. % TODO: Document specific frontend framework (e.g., React, Vue, plain HTML/JS).
    \item \textbf{Backend:} A server-side application that orchestrates data ingestion, preprocessing, inference, post-processing, and output generation. The backend exposes an API that the frontend calls to submit volumes and retrieve results. % TODO: Document backend framework (e.g., FastAPI, Flask, Django).
    \item \textbf{Inference engine:} The 3D U-Net model and associated pipeline (preprocessing, sliding-window inference, TTA, post-processing) are invoked by the backend to process uploaded volumes.
\end{itemize}

Communication between frontend and backend occurs over HTTP/HTTPS. The 3D visualization is rendered client-side using a WebGL-based library. % TODO: Document 3D visualization library (e.g., Three.js, VTK.js, vtk.js).

The system is designed to run on a single server with GPU acceleration for inference. No distributed computing, load balancing, or multi-user concurrency support is implemented in the current prototype.

\subsection{Data Ingestion}

The system accepts CT volumes in NIfTI format (\texttt{.nii} or \texttt{.nii.gz}), the standard format for medical imaging data that preserves spatial metadata including voxel dimensions, spacing, and orientation.

\begin{itemize}
    \item \textbf{Upload:} Users upload a NIfTI file through the web interface. The frontend transmits the file to the backend via HTTP multipart form data.
    \item \textbf{Validation:} The backend validates the uploaded file for correct format, readable metadata, and non-corrupt data. Files that fail validation are rejected with an error message. % TODO: Document specific validation checks (file size limits, format verification, metadata consistency).
    \item \textbf{Storage:} Uploaded files are stored temporarily on the server for processing. % TODO: Document whether uploaded files are deleted after processing or retained. If retained, document retention policy and access controls.
\end{itemize}

The ingestion pipeline is designed to handle single-volume uploads. Batch processing of multiple volumes is not implemented in the current prototype.

\subsection{Backend Processing}

Once a volume is uploaded and validated, the backend executes the segmentation pipeline:

\begin{enumerate}
    \item \textbf{Preprocessing:} The uploaded NIfTI volume is loaded, validated for metadata consistency, clipped to $[-200, 300]$ HU, and normalized to $[0, 1]$ using min--max scaling. This preprocessing is identical to that applied during training.
    \item \textbf{Inference:} The preprocessed volume is processed through the 3D U-Net using sliding-window inference with $64 \times 192 \times 192$ voxel patches, 50\% overlap, and 8-flip test-time augmentation. The inference engine runs on GPU hardware with sufficient memory to process patches sequentially.
    \item \textbf{Post-processing:} The raw segmentation mask is refined using connected-component filtering (kidney blobs $< 5000$ voxels removed, tumor blobs $< 100$ voxels removed) to eliminate spurious detections.
\end{enumerate}

The backend processing pipeline is deterministic and reproducible. Identical preprocessing, inference, and post-processing are applied to all uploaded volumes.

\subsection{Output Generation}

The segmentation pipeline produces three types of output:

\begin{itemize}
    \item \textbf{Segmentation mask:} A 3D volume of the same spatial dimensions as the input, where each voxel is assigned one of three class labels (background, kidney, tumor/cyst). The mask is stored in NIfTI format and can be downloaded by the user.
    \item \textbf{Volumetric statistics:} Kidney volume (ml), tumor/cyst volume (ml), and tumor detection status (binary indicator) are computed from the segmentation mask using voxel spacing metadata. These statistics are formatted for display in the web interface.
    \item \textbf{3D mesh visualization:} The kidney and tumor/cyst segmentation masks are converted to 3D surface meshes using the marching cubes algorithm. The meshes are exported in a format compatible with WebGL rendering (e.g., STL, OBJ, or custom binary format). % TODO: Document exact mesh export format.
\end{itemize}

All outputs are generated server-side and transmitted to the frontend for display. The segmentation mask and volumetric statistics can be downloaded for offline analysis.

\subsection{Web Interface}

The browser-based web interface provides an accessible visualization of segmentation results without requiring specialized medical imaging software:

\begin{itemize}
    \item \textbf{Upload panel:} A file input widget allows users to select and upload a NIfTI CT volume. Progress indicators display upload status and processing status. % TODO: Document whether progress indicators are implemented for long-running inference.
    \item \textbf{3D viewer:} An interactive WebGL-based viewer displays the kidney and tumor/cyst segmentation meshes. Users can rotate, zoom, and pan the 3D view to inspect segmentation results from any angle. Individual classes (kidney, tumor/cyst) can be toggled on or off.
    \item \textbf{Volumetric statistics panel:} Kidney volume, tumor/cyst volume, and tumor detection status are displayed alongside the 3D visualization in a text-based panel.
    \item \textbf{Download links:} Users can download the segmentation mask (NIfTI format) and volumetric statistics (text or JSON format) for offline analysis.
\end{itemize}

The web interface is designed for single-user interaction. Multi-user access, authentication, and session management are not implemented in the current prototype.

\begin{figure}[!t]
\centering
% TODO: Replace with actual figure: fig_webapp.pdf
% \includegraphics[width=0.9\linewidth]{figures/fig_webapp.pdf}
\fbox{\parbox{0.7\linewidth}{\centering\vspace{2cm}\textbf{[Placeholder]}\\Web Application Screenshot\\\small Upload panel | 3D Viewer | Volumetric Statistics\vspace{2cm}}}
\caption{Browser-based web application interface. Users upload a NIfTI CT volume, and the system displays an interactive 3D segmentation mesh alongside kidney and tumor volumetric statistics.}
\label{fig:webapp}
\end{figure}

\subsection{Safety UX}

The web interface includes explicit safety messaging to prevent misuse and ensure that users understand the system's limitations:

\begin{itemize}
    \item \textbf{Disclaimer:} A prominent disclaimer is displayed on the interface stating that NephroVision is an academic prototype for decision-support use and is not clinically approved.
    \item \textbf{Physician review message:} All segmentation outputs are accompanied by a message stating that results require physician review before any clinical application.
    \item \textbf{Decision-support wording:} The interface consistently uses decision-support language (e.g., ``segmentation support,'' ``volumetric statistics for review'') rather than diagnostic language (e.g., ``diagnosis,'' ``detection of cancer'').
    \item \textbf{Limitations statement:} The interface includes a brief statement of known limitations, including moderate tumor boundary precision and the merged tumor/cyst class convention.
\end{itemize}

These safety messages are designed to prevent over-reliance on automated outputs and to reinforce the decision-support role of the system. The disclaimer and physician review message are displayed prominently to ensure visibility.

\subsection{Privacy and Security Considerations}

The current prototype implements basic privacy and security measures, but formal cybersecurity validation has not been conducted:

\begin{itemize}
    \item \textbf{Uploaded file retention:} % TODO: Document whether uploaded files are deleted after processing. If files are retained, document retention policy, access controls, and encryption at rest.
    \item \textbf{Model weights:} The trained 3D U-Net model weights are stored on the server and are not exposed to users. % TODO: Document whether model weights are accessible via the API or protected by access controls.
    \item \textbf{Network security:} Communication between frontend and backend occurs over HTTP. % TODO: Document whether HTTPS/TLS is implemented for encrypted communication.
    \item \textbf{Input validation:} Uploaded files are validated for format and metadata consistency. Malformed or corrupted files are rejected. % TODO: Document additional input validation measures (e.g., file size limits, content-type verification).
    \item \textbf{Cybersecurity hardening:} Formal cybersecurity assessment, penetration testing, and compliance with medical device cybersecurity standards (e.g., IEC 81001-5-1) are identified as future work.
\end{itemize}

The system is not designed for production deployment in clinical environments. Privacy and security measures are appropriate for academic prototype use but require formal validation before any clinical deployment.

\subsection{Maintainability}

The software system is designed to support maintainability and reproducibility:

\begin{itemize}
    \item \textbf{Configuration files:} Preprocessing parameters (HU clipping range, normalization), inference configuration (patch size, overlap, TTA), and post-processing thresholds (blob removal) are stored in configuration files rather than hardcoded. This enables modification without code changes.
    \item \textbf{Modular code:} The system is organized into modular components (data ingestion, preprocessing, inference, post-processing, visualization) with clear interfaces. This modularity supports independent testing and future extension.
    \item \textbf{Reproducibility documentation:} The project maintains comprehensive documentation of environment specifications, random seeds, commands, and checkpoint information in \texttt{reproducibility\_notes.md}. This documentation enables reproduction of the segmentation pipeline on new hardware.
    \item \textbf{Version control:} The codebase is maintained in version control with commit history documenting changes. % TODO: Document version control system (e.g., Git) and repository structure.
\end{itemize}

The maintainability design supports future development, including extension to multi-center data, integration with clinical workflows, and formal cybersecurity hardening.
